% Luc Destombe --- licence CC BY-NC-SA 4.0
% https://creativecommons.org/licenses/by-nc-sa/4.0/deed.fr
% Fichier autonome : compiler avec pdflatex, deux fois (signets, nombre de pages).
\documentclass[10pt,a4paper,twocolumn]{article}
\makeatletter
\newif\iflycee@unecolonne
\newif\iflycee@palatino
\lycee@palatinotrue

\RequirePackage[utf8]{inputenc}
\RequirePackage[T1]{fontenc}
\RequirePackage[french]{babel}
\RequirePackage{etoolbox}

\newcommand{\ShorthandsOff}{\@ifundefined{shorthandoff}{}{\shorthandoff{;:!?}}}
\newcommand{\ShorthandsOn}{\@ifundefined{shorthandon}{}{\shorthandon{;:!?}}}
\ShorthandsOff
\AtBeginDocument{\ShorthandsOn}
\AtBeginEnvironment{tikzpicture}{\ShorthandsOff}

\RequirePackage{geometry}
\RequirePackage{fancyhdr}
\RequirePackage{lastpage}
\RequirePackage{multicol}
\RequirePackage{array}
\RequirePackage{enumitem}
\RequirePackage{xcolor}
\RequirePackage{needspace}

\RequirePackage{amsmath,amssymb}
\RequirePackage{mathpazo}
\DeclareMathSizes{10.95}{10.95}{8}{5.5}
\begingroup\catcode`\_=\active \catcode`\^=\active
\gdef_#1{\sb{\scriptscriptstyle #1}}
\gdef^#1{\sp{\scriptscriptstyle\lycee@moinsepais #1}}
\endgroup
\newcommand\lycee@moins{\mathbin{%
  \setbox\z@\hbox{$\m@th\scriptscriptstyle\lycee@moinsorig$}%
  \dimen@=.55\wd\z@ \advance\dimen@-.5pt
  \hbox to.75\wd\z@{\hss$\m@th\scriptscriptstyle\vcenter{\hbox{%
    \tikz\draw[line width=.5pt,line cap=round](0,0)--(\the\dimen@,0);}}$\hss}}}
\begingroup\lccode`\~=`\- \lowercase{\endgroup\let~}\lycee@moins
\newcommand\lycee@moinsepais{\mathcode`\-="8000 }
\AtBeginDocument{\mathchardef\lycee@moinsorig=\mathcode`\-\relax}
\AtBeginDocument{\catcode`\_=12 \mathcode`\_="8000
  \catcode`\^=12 \mathcode`\^="8000 }
\newcommand\lycee@exposants{%
  \fontdimen13\textfont2=.48\fontdimen6\textfont2
  \fontdimen14\textfont2=.48\fontdimen6\textfont2
  \fontdimen15\textfont2=.48\fontdimen6\textfont2
  \fontdimen13\scriptfont2=.48\fontdimen6\scriptfont2
  \fontdimen14\scriptfont2=.48\fontdimen6\scriptfont2
  \fontdimen15\scriptfont2=.48\fontdimen6\scriptfont2}
\every@math@size\expandafter{\the\every@math@size\lycee@exposants}
\newlength{\retraitcalculs}
\setlength{\retraitcalculs}{2em}

\RequirePackage{tikz}
\usetikzlibrary{arrows.meta,calc,babel,shapes.misc}
\RequirePackage[most]{tcolorbox}
\tcbuselibrary{skins,breakable}

\definecolor{coulDef}{HTML}{1F4E79}
\definecolor{coulProp}{HTML}{7030A0}
\definecolor{coulRem}{HTML}{C55A11}
\definecolor{coulEx}{HTML}{2E7D32}
\definecolor{coulExo}{HTML}{B03A2E}
\definecolor{coulPart}{HTML}{1F4E79}
\definecolor{coulMeth}{HTML}{1B6E4F}
\definecolor{coulCourbe}{HTML}{0B5ED7}
\definecolor{coulCourbeB}{HTML}{C00000}
\definecolor{coulCourbeC}{HTML}{2E7D32}
\definecolor{coulGrille}{HTML}{8C8C8C}
\definecolor{coulRep}{HTML}{1B6E4F}
\definecolor{coulBar}{HTML}{2E7D32}

\newcommand{\R}{\mathbb{R}}
\newcommand{\Rs}{\mathbb{R}^{*}}
\renewcommand{\le}{\leqslant}
\renewcommand{\ge}{\geqslant}
\newcommand{\vect}[1]{\overrightarrow{#1}}
\newcommand{\norme}[1]{\left\lVert #1 \right\rVert}
\newcommand{\intervalle}[2]{\left[#1\,;#2\right]}
\RequirePackage{cancel}
\renewcommand{\CancelColor}{\color{coulExo}}
\newcommand{\fois}[1]{\times{\color{coulExo}#1}}
\newcommand{\barre}[1]{\cancel{#1}}

\tikzset{grille/.style={coulGrille,line width=0.4pt}}
\newcommand{\quadrillage}[4]{\draw[grille,step=1] (#1,#2) grid (#3,#4);}
\tikzset{
  croix/.style={cross out,draw,line width=0.6pt,inner sep=0pt,minimum size=4pt},
  croixdroite/.style={croix,rotate=45,line width=0.8pt},
}
\newcommand{\pt}[3]{\path (#1) node[croix]{} node[#2,font=\small,inner sep=2.5pt]{$#3$};}

\newcommand{\lycee@repere}[4]{%
  \draw[grille,step=1] (#1,#2) grid (#3,#4);
  \draw[-{Stealth[length=2mm]},thick] (#1,0)--({#3+0.4},0) node[below left=-1pt]{$x$};
  \draw[-{Stealth[length=2mm]},thick] (0,#2)--(0,{#4+0.4}) node[below left=-1pt]{$y$};
  \node[below left,font=\scriptsize,inner sep=1.5pt] at (0,0) {$0$};}
\newcommand{\lycee@gradx}[1]{\foreach \k/\l in {#1}{%
  \draw (\k,0.07)--(\k,-0.07) node[below,font=\scriptsize,inner sep=1.5pt]{$\l$};}}
\newcommand{\lycee@grady}[1]{\foreach \k/\l in {#1}{%
  \draw (0.07,\k)--(-0.07,\k) node[left,font=\scriptsize,inner sep=1.5pt]{$\l$};}}
\newcommand{\lycee@intff}[2]{\left[#1\,;#2\right]}
\newcommand{\lycee@intfo}[2]{\left[#1\,;#2\right[}
\newcommand{\lycee@intof}[2]{\left]#1\,;#2\right]}
\newcommand{\lycee@intoo}[2]{\left]#1\,;#2\right[}
\newcommand{\lycee@ens}[1]{\left\{#1\right\}}
\AtBeginDocument{%
  \providecommand{\repere}{\lycee@repere}%
  \providecommand{\gradx}{\lycee@gradx}%
  \providecommand{\grady}{\lycee@grady}%
  \providecommand{\Cf}{\mathcal{C}\sb{\scriptscriptstyle f}}%
  \providecommand{\Cg}{\mathcal{C}\sb{\scriptscriptstyle g}}%
  \providecommand{\intff}{\lycee@intff}%
  \providecommand{\intfo}{\lycee@intfo}%
  \providecommand{\intof}{\lycee@intof}%
  \providecommand{\intoo}{\lycee@intoo}%
  \providecommand{\ens}{\lycee@ens}}

\newlist{questions}{enumerate}{3}
\setlist[questions,1]{label=\arabic*\textdegree),leftmargin=*,itemsep=2pt,topsep=3pt}
\setlist[questions,2]{label=\alph*),leftmargin=*,itemsep=1pt,topsep=2pt}
\setlist[questions,3]{label=\roman*),leftmargin=*,itemsep=1pt,topsep=1pt}
\setlist[itemize]{label=\textbullet,leftmargin=*,itemsep=2pt,topsep=3pt}

\newcommand{\besoinplace}[1]{\par\penalty\@M\begingroup
  \dimen@=#1\relax\dimen@ii=\pagegoal\advance\dimen@ii-\pagetotal
  \ifdim\dimen@>\dimen@ii\ifdim\dimen@ii>\z@\vfil\fi\break\fi\endgroup}

\newcommand{\lycee@pied}{\thepage/\pageref{LastPage}}
\newcommand{\entete}[2]{%
  \pagestyle{fancy}\fancyhf{}%
  \fancyhead[L]{\footnotesize\color{coulPart}\textbf{#1}}%
  \fancyhead[R]{\footnotesize\color{coulPart}\textbf{#2}}%
  \fancyfoot[C]{\footnotesize\lycee@pied}%
  \renewcommand{\headrulewidth}{0.4pt}%
  \renewcommand{\headrule}{\hbox to\headwidth{%
    \color{coulPart!40}\leaders\hrule height \headrulewidth\hfill}}}

\RequirePackage{ccicons}
\newcommand{\lycee@licence}{{\scriptsize\color{gray}%
  \copyright\ Luc Destombe --- \ccbyncsa\ CC BY-NC-SA 4.0}}
\newif\iflycee@licence \lycee@licencetrue
\newcommand{\sanslicence}{\lycee@licencefalse}
\AtBeginDocument{\iflycee@licence\AddToHookNext{shipout/foreground}{%
  \put(0,0){\rlap{\kern\dimexpr1in+\hoffset+\oddsidemargin+\textwidth\relax
    \llap{\raisebox{-\dimexpr1in+\voffset+\topmargin+\headheight+\headsep
      +\textheight+\footskip\relax}{\lycee@licence}}}}}\fi}

\newcommand{\formulesserrees}{%
  \setlength{\abovedisplayskip}{3pt}\setlength{\belowdisplayskip}{3pt}%
  \setlength{\abovedisplayshortskip}{2pt}\setlength{\belowdisplayshortskip}{3pt}}

\setlength{\parindent}{0pt}

\RequirePackage[implicit=false,hidelinks,pdfencoding=auto,psdextra,bookmarksopen,
  bookmarksopenlevel=1,pdfcreator={},pdfproducer={}]{hyperref}
\RequirePackage{bookmark}
\hypersetup{pdfauthor={Luc Destombe},pdfkeywords={CC BY-NC-SA 4.0}}
\newcounter{lycee@signet}
\newcommand{\signet}[3][]{\stepcounter{lycee@signet}%
  \edef\lycee@dest{\if\relax\detokenize{#1}\relax
    signet.\arabic{lycee@signet}\else#1\fi}%
  \hypertarget{\lycee@dest}{}\bookmark[level=#2,dest=\lycee@dest]{#3}}
\pdfstringdefDisableCommands{%
  \def\R{R}\def\vect#1{#1}\def\Cf{Cf}\def\Cg{Cg}\def\fois#1{×#1}%
  \def\barre#1{#1}\def\intervalle#1#2{[#1;#2]}\def\intff#1#2{[#1;#2]}%
  \def\textsuperscript#1{#1}\def\dfrac#1#2{#1/#2}\def\frac#1#2{#1/#2}%
  \def\vec#1{vecteur #1}}
\RequirePackage[official]{eurosym}

\tcbset{exercice corrige/.style={
  enhanced,breakable,before upper=\raggedright,
  colback=white,colframe=coulExo,
  boxrule=0.8pt,arc=2pt,
  left=6pt,right=6pt,top=8pt,bottom=5pt,
  before skip=9pt,after skip=4pt,
  coltitle=white,fonttitle=\bfseries\small,
  attach boxed title to top left={xshift=8pt,yshift=-\tcboxedtitleheight/2},
  boxed title style={colback=coulExo,arc=2pt,boxrule=0pt}}}
\newcounter{exo}
\newtcolorbox{exercice}[1][]{exercice corrige,
  phantom={\signet[exercice-\arabic{exo}]{2}{Exercice \arabic{exo}}},
  title={Exercice \theexo\ifx\relax#1\relax\else{} --- #1\fi}}
\newenvironment{exo}[1][\relax]{\refstepcounter{exo}\begin{exercice}[#1]}{\end{exercice}}

\RequirePackage{cuted}
\geometry{top=1.6cm,bottom=1cm,left=1cm,right=1cm,headheight=14pt,footskip=0.6cm}

\newcommand{\entetecorrige}[3]{%
  \begin{strip}
  \begin{center}
    \begin{tcolorbox}[enhanced,width=0.92\linewidth,colback=white,colframe=coulPart,
        boxrule=1.6pt,arc=3pt,halign=center,top=5pt,bottom=5pt]
      {\LARGE\bfseries\color{coulPart} #1}\\[3pt]
      {\normalsize #2}
    \end{tcolorbox}
  \end{center}
  \if\relax\detokenize{#3}\relax\else

  \vspace{2pt}
  \noindent
  \begin{tcolorbox}[enhanced,colback=white,colframe=black!35,boxrule=0.5pt,arc=2pt,
      left=7pt,right=7pt,top=4pt,bottom=4pt]
  {\small #3}
  \end{tcolorbox}
  \vspace{6pt}
  \fi
  \end{strip}}

\newcounter{partie}
\newcommand{\partiebandeau}[1]{%
  \besoinplace{8\baselineskip}\vspace{6pt}%
  \begin{tcolorbox}[enhanced,colback=coulPart!8,colframe=coulPart,boxrule=0pt,
      leftrule=3pt,arc=0pt,left=5pt,right=4pt,top=3pt,bottom=3pt,
      before skip=4pt,after skip=2pt]
    \bfseries\color{coulPart}#1\signet{1}{#1}
  \end{tcolorbox}}
\newcommand{\partie}[1]{\stepcounter{partie}\partiebandeau{\Roman{partie}) #1}}
\newcommand{\partiesansnum}[1]{\partiebandeau{#1}}

\newcommand{\res}[1]{%
  \tcbox[on line,colback=coulPart!7,colframe=coulPart,boxrule=0.5pt,arc=1pt,
    boxsep=0pt,left=3pt,right=3pt,top=1pt,bottom=2pt]{$#1$}}
\newcommand{\calc}[1]{\par\smallskip\noindent$\displaystyle #1$\par}
\newenvironment{calculs}{\par\smallskip\noindent$\displaystyle\begin{aligned}[t]}%
  {\end{aligned}$\par\smallskip}
\newcommand{\rem}[1]{\par\smallskip{\small\itshape\color{coulPart}#1}\par}
\newcommand{\deux}[2]{\par\noindent\makebox[0.5\linewidth][l]{#1}#2}

\setlist[questions,1]{label=\arabic*\textdegree),leftmargin=*,itemsep=2pt,topsep=2pt}
\setlist[questions,2]{label=\alph*),leftmargin=*,itemsep=2pt,topsep=1pt}
\newlist{reponses}{enumerate}{1}
\setlist[reponses]{label=\alph*),leftmargin=*,itemsep=2pt,topsep=2pt}

\setlength{\parskip}{1pt}
\setlength{\columnsep}{0.6cm}
\raggedbottom
\formulesserrees
\AtBeginDocument{\formulesserrees}

\makeatother
\entete{Chapitre 1 --- Calcul littéral --- Corrigé des exercices}{1\textsuperscript{re} STI2D}

\usepackage{tkz-tab}
\let\deuxpaquet\deux
\renewcommand{\deux}[2]{\deuxpaquet{#1}{#2}\par}
\newenvironment{tabsignes}[3][2.2]{\begin{center}\begin{tikzpicture}[scale=0.62]
  \tkzTabInit[lgt=#1,espcl=1.7]{#2}{#3}}{\end{tikzpicture}\end{center}}
\newcommand{\petittab}[3]{\begin{minipage}{0.49\linewidth}\centering
  \begin{tikzpicture}[scale=0.55]
  \tkzTabInit[lgt=2.9,espcl=1.2]{$x$/1,#1/1}{$-\infty$,#2,$+\infty$}
  \tkzTabLine{,#3,}
  \end{tikzpicture}\end{minipage}}
\clubpenalty=10000
\widowpenalty=10000

\begin{document}

\entetecorrige{CORRIGÉ --- CALCUL LITTÉRAL}
  {Chapitre 1 : fiche d'exercices --- Première STI2D}
  {Les résultats sont encadrés ; les remarques en bleu signalent les erreurs
   fréquentes et les méthodes à retenir. À partir de la partie V, chaque équation ou
   inéquation se termine par l'ensemble $S$ des solutions.}

\partie{Réduire une expression littérale}

\begin{exo}[Réduire une somme]
\rem{On ne regroupe que des termes semblables : les $a$ entre eux, les $a^{2}$ entre
eux, les nombres entre eux.}
\textbf{a)}\par
\deux{$A=\res{11a+3}$}{$B=\res{16a}$}
\deux{$C=\res{10a+6}$}{$D=\res{15a^{2}}$}
\deux{$E=\res{6a^{2}+9}$}{$F=\res{8a^{2}+3a}$}
\textbf{b)}\par
\deux{$G=\res{-2m}$}{$H=\res{m}$}
\deux{$I=\res{-15m}$}{$J=\res{4m^{2}}$}
\deux{$K=\res{-5m^{2}+5m}$}{$L=\res{-7m^{2}}$}
\textbf{c)}
\begin{calculs}
M&=(3n^{2}+5n^{2})+4n+(-6+1+8)\\
M&=\res{8n^{2}+4n+3}\\
N&=(-7n^{2}+3n^{2})+(-5n-n)+(9-4)\\
N&=\res{-4n^{2}-6n+5}\\
O&=(-6n^{2}-2n^{2})+(5n-9n)+(-3+8)\\
O&=\res{-8n^{2}-4n+5}\\
P&=(-3n^{2}+7n^{2})+(6n+2n)+(4-4)\\
P&=\res{4n^{2}+8n}
\end{calculs}
\rem{Dans $C=6+2a+8a$, le $6$ reste seul : on n'écrit pas $16a$.}
\end{exo}

\begin{exo}[Réduire un produit]
\rem{On multiplie les nombres entre eux (règle des signes), puis les lettres entre
elles : $x\times x=x^{2}$.}
\textbf{a)}\par
\deux{$A=\res{16x}$}{$B=\res{12x^{2}}$}
\deux{$C=\res{10x^{2}}$}{$D=\res{21x}$}
\deux{$E=\res{18x^{2}}$}{$F=\res{20x^{2}}$}
\textbf{b)}\par
\deux{$G=\res{-12x^{2}}$}{$H=\res{-27x}$}
\deux{$I=\res{-10x^{2}}$}{$J=\res{21x^{2}}$}
\deux{$K=\res{-12x^{2}}$}{$L=\res{-15x^{2}}$}
\deux{$M=\res{12x^{2}}$}{$N=\res{-12x^{2}}$}
\deux{$O=\res{-14x^{2}}$}{}
\end{exo}

\begin{exo}[Somme ou produit ?]
\deux{$A=5+4x$ : \res{\text{ne se réduit pas}}}{}
\deux{$B=\res{20x}$}{$C=\res{20x^{2}}$}
\deux{$D=\res{9x}$}{}
\rem{$A$ est une somme d'un nombre et d'un terme en $x$, qui ne sont pas semblables.
$B$ et $C$ sont des produits : ils se réduisent toujours. $D$ est une somme de deux
termes semblables : $5x+4x=9x$, et non $20x^{2}$.}
\end{exo}

\partie{Développer}

\begin{exo}[Distributivité simple]
\textbf{a)}\par
\deux{$A=\res{12x+20}$}{$B=\res{12+10x}$}
\deux{$C=\res{2x^{2}+9x}$}{$D=\res{12x^{2}+15x}$}
\textbf{b)}
\begin{calculs}
E&=5x+6x+12\\
E&=\res{11x+12}
\end{calculs}
\deux{$F=\res{18x+24}$}{$G=\res{17x+18}$}
\deux{$H=\res{18x+17}$}{$I=\res{3x^{2}+17x+2}$}
\begin{calculs}
J&=8x^{2}+6x+6x^{2}+15x\\
J&=\res{14x^{2}+21x}
\end{calculs}
\rem{Dans $E$, on développe d'abord $3(2x+4)$ : on n'additionne pas $5x$ et $3$.}
\end{exo}

\begin{exo}[Signes et facteurs négatifs]
\rem{Le signe devant le facteur fait partie du facteur : $-3(5+2a)=-15-6a$.}
\textbf{a)}
\begin{calculs}
A&=5a+15-4a & B&=-7a+6a-12\\
A&=\res{a+15} & B&=\res{-a-12}\\
C&=8+12a-15-6a & D&=15a-10+8a-6\\
C&=\res{6a-7} & D&=\res{23a-16}
\end{calculs}
\textbf{b)}
\begin{calculs}
E&=3b+10b+12b^{2}\\
E&=\res{12b^{2}+13b}\\
F&=-12b^{2}+20b+7b^{2}\\
F&=\res{-5b^{2}+20b}\\
G&=-12b^{2}-6b-2b^{2}+9b\\
G&=\res{-14b^{2}+3b}\\
H&=-10b^{2}+6b-12b^{2}-24b\\
H&=\res{-22b^{2}-18b}
\end{calculs}
\rem{Dans $G$ : $-b(2b-9)=-2b^{2}+9b$ ; le $-$ change les deux signes.}
\end{exo}

\begin{exo}[Double distributivité]
\rem{Chaque terme de la première parenthèse multiplie chaque terme de la seconde :
quatre produits, puis on réduit.}
\begin{calculs}
A&=8x^{2}+6x+20x+15\\
A&=\res{8x^{2}+26x+15}\\
B&=15x+6+20x^{2}+8x\\
B&=\res{20x^{2}+23x+6}\\
C&=20x+15x^{2}+8+6x\\
C&=\res{15x^{2}+26x+8}\\
D&=18x+30x^{2}+3+5x\\
D&=\res{30x^{2}+23x+3}\\
E&=6x+10+15x\\
E&=\res{21x+10}\\
F&=16x^{2}+12x+12x+9\\
F&=\res{16x^{2}+24x+9}
\end{calculs}
\rem{$E$ n'est pas un produit de deux parenthèses : seul $5$ multiplie $(2+3x)$,
c'est une distributivité simple. $F=(4x+3)^{2}$ : une identité remarquable ira plus
vite (partie III).}
\end{exo}

\partie{Les trois identités remarquables}

\begin{exo}[Identité $(a+b)^{2}$]
\rem{$(a+b)^{2}=a^{2}+2ab+b^{2}$ : ne pas oublier le double produit $2ab$.}
\begin{calculs}
A&=(3x)^{2}+2\times 3x\times 2+2^{2}\\
A&=\res{9x^{2}+12x+4}
\end{calculs}
\deux{$B=\res{9x^{2}+24x+16}$}{$C=\res{4a^{2}+28a+49}$}
\deux{$D=\res{16b^{2}+40b+25}$}{$E=\res{x^{2}+18x+81}$}
\deux{$F=\res{36x^{2}+12x+1}$}{$G=\res{4m^{2}+20m+25}$}
\deux{$H=\res{49r^{2}+42r+9}$}{}
\rem{$(3x)^{2}=9x^{2}$ et non $3x^{2}$ : le carré porte sur le nombre \emph{et} sur la
lettre.}
\end{exo}

\begin{exo}[Identité $(a-b)^{2}$]
\begin{calculs}
A&=(2x)^{2}-2\times 2x\times 3+3^{2}\\
A&=\res{4x^{2}-12x+9}
\end{calculs}
\deux{$B=\res{25x^{2}-40x+16}$}{$C=\res{16-24x+9x^{2}}$}
\deux{$D=\res{1-12x+36x^{2}}$}{$E=\res{9a^{2}-30a+25}$}
\deux{$F=\res{4-28b+49b^{2}}$}{$G=\res{16m^{2}-8m+1}$}
\deux{$H=\res{n^{2}-16n+64}$}{}
\rem{Le dernier terme $b^{2}$ est toujours positif : seul le double produit porte le
signe $-$.}
\end{exo}

\begin{exo}[Identité $(a+b)(a-b)$]
\begin{calculs}
A&=(3x)^{2}-7^{2}\\
A&=\res{9x^{2}-49}\\
B&=4x^{2}-10x+6x-15\\
B&=\res{4x^{2}-4x-15}
\end{calculs}
\deux{$C=\res{25-4x^{2}}$}{$D=\res{16a^{2}-1}$}
\deux{$E=\res{b^{2}-36}$}{$F=\res{81-c^{2}}$}
\rem{$B$ n'est pas une identité remarquable : $+3$ et $-5$ ne sont pas opposés ; on
utilise la double distributivité.}
\end{exo}

\begin{exo}[Choisir la bonne méthode]
\textbf{a)}\par
\deux{$A=\res{6x+21}$}{$B=\res{4x^{2}+28x+49}$}
\deux{$C=\res{25x^{2}-30x+9}$}{$D=\res{49x^{2}-1}$}
\begin{calculs}
E&=12x^{2}-18x+10x-15\\
E&=\res{12x^{2}-8x-15}
\end{calculs}
\textbf{b)}
\begin{calculs}
F&=3x^{2}+12x+2x+8-\big(2x^{2}+2x-5x-5\big)\\
F&=3x^{2}+14x+8-\big(2x^{2}-3x-5\big)\\
F&=3x^{2}+14x+8-2x^{2}+3x+5\\
F&=\res{x^{2}+17x+13}
\end{calculs}
\begin{calculs}
G&=x^{2}-6x+9+4x^{2}-1\\
G&=\res{5x^{2}-6x+8}
\end{calculs}
\begin{calculs}
H&=4x^{2}+20x+25-\big(3x^{2}+2x-12x-8\big)\\
H&=4x^{2}+20x+25-\big(3x^{2}-10x-8\big)\\
H&=4x^{2}+20x+25-3x^{2}+10x+8\\
H&=\res{x^{2}+30x+33}
\end{calculs}
\rem{Le signe $-$ devant un produit change \emph{tous} les signes : on garde la
parenthèse tant que le produit n'est pas réduit. C'est l'erreur n°1 du chapitre.}
\end{exo}

\partie{Factoriser}

\begin{exo}[Facteur commun]
\rem{On met en facteur le plus grand diviseur commun des coefficients, et la plus
grande puissance de $x$ présente dans chaque terme.}
\deux{$A=\res{6(a+b)}$}{$B=\res{12(4a-3b)}$}
\deux{$C=\res{15(3a-b)}$}{$D=\res{40(3a-2b)}$}
\deux{$E=\res{x(7x+4)}$}{$F=\res{x(18x-11)}$}
\deux{$G=\res{x(15+8x)}$}{$H=\res{2x(5x-2)}$}
\deux{$I=\res{x(4x+1)}$}{$J=\res{x^{2}(3x-1)}$}
\rem{Dans $I$, $x=x\times 1$ : il reste $1$ dans la parenthèse, pas $0$. On vérifie
en redéveloppant.}
\end{exo}

\begin{exo}[Identité $a^{2}+2ab+b^{2}$]
\rem{On repère $a^{2}$ et $b^{2}$, puis on \emph{vérifie} que le troisième terme est
bien $2ab$.}
\begin{calculs}
A&=(3x)^{2}+2\times 3x\times 2+2^{2}\\
A&=\res{(3x+2)^{2}}\\
B&=(5x)^{2}+2\times 5x\times 4+4^{2}\\
B&=\res{(5x+4)^{2}}
\end{calculs}
$C$ : $(3a)^{2}$ et $2^{2}$, mais $2\times 3a\times 2=12a$, pas $6a$ :
\res{\text{pas un carré}}
\begin{calculs}
D&=5^{2}+2\times 5\times 3b+(3b)^{2}\\
D&=\res{(5+3b)^{2}}
\end{calculs}
$E$ : $(4c)^{2}$ et $5^{2}$, mais $2\times 4c\times 5=40c$, pas $20c$ :
\res{\text{pas un carré}}
\begin{calculs}
F&=(7d)^{2}+2\times 7d\times 2+2^{2}\\
F&=\res{(7d+2)^{2}}
\end{calculs}
\rem{Pour $F$, on range d'abord : $49d^{2}+28d+4$.}
\end{exo}

\begin{exo}[Identités $a^{2}-2ab+b^{2}$ et $a^{2}-b^{2}$]
\textbf{a)}
\begin{calculs}
A&=(4x)^{2}-2\times 4x\times 3+3^{2}\\
A&=\res{(4x-3)^{2}}\\
B&=3^{2}-2\times 3\times 2x+(2x)^{2}\\
B&=\res{(3-2x)^{2}}
\end{calculs}
$C$ : $2\times 2x\times 3=12x$, pas $6x$ : \res{\text{pas un carré}}\par
$D$ : le terme $-1$ est négatif : \res{\text{pas un carré}}
\begin{calculs}
E&=(3a)^{2}-2\times 3a\times 5+5^{2}\\
E&=\res{(3a-5)^{2}}\\
F&=(7b)^{2}-2\times 7b\times 2+2^{2}\\
F&=\res{(7b-2)^{2}}
\end{calculs}
\textbf{b)}
\begin{calculs}
G&=(7x)^{2}-5^{2}\\
G&=\res{(7x-5)(7x+5)}\\
H&=8^{2}-(3x)^{2}\\
H&=\res{(8-3x)(8+3x)}\\
I&=\big[(x+2)-4\big]\big[(x+2)+4\big]\\
I&=\res{(x-2)(x+6)}\\
J&=\big[6-(x-1)\big]\big[6+(x-1)\big]\\
J&=\res{(7-x)(5+x)}
\end{calculs}
\rem{Un carré est toujours positif : si l'un des deux termes « carrés » est négatif
($D$), ce n'est pas une identité remarquable. Dans $J$, les crochets évitent l'erreur
de signe : $6-(x-1)=7-x$.}
\end{exo}

\begin{exo}[Facteur commun entre parenthèses]
\rem{On écrit le facteur commun devant, puis entre crochets ce qui reste de chaque
terme ; on réduit le crochet.}
\begin{calculs}
A&=(3x+4)\big[(2x-5)+(5x+1)\big]\\
A&=\res{(3x+4)(7x-4)}\\
B&=(6x-1)\big[(2x+3)+(x-5)\big]\\
B&=\res{(6x-1)(3x-2)}\\
C&=(4x-1)\big[(2x+9)+(3x+2)\big]\\
C&=\res{(4x-1)(5x+11)}\\
D&=(2r+5)\big[(4r-3)+(r+7)\big]\\
D&=\res{(2r+5)(5r+4)}\\
E&=(5x+3)(x-4)+(5x+3)\times 1\\
E&=(5x+3)\big[(x-4)+1\big]\\
E&=\res{(5x+3)(x-3)}\\
F&=(7r-2)(3r+1)+(7r-2)(7r-2)\\
F&=(7r-2)\big[(3r+1)+(7r-2)\big]\\
F&=\res{(7r-2)(10r-1)}\\
G&=(4m-5)(4m-5)-(4m-5)(m+3)\\
G&=(4m-5)\big[(4m-5)-(m+3)\big]\\
G&=(4m-5)(4m-5-m-3)\\
G&=\res{(4m-5)(3m-8)}
\end{calculs}
\rem{Dans $G$, le $-$ devant le second produit change les deux signes de $(m+3)$.
Dans $E$, sans le $\times 1$, on oublie le $+1$.}
\end{exo}

\partie{Résoudre des équations}

\begin{exo}[Équations du premier degré]
\begin{reponses}
  \item $4x-3x=12+9$, donc $x=21$ : \ \res{S=\ens{21}}
  \item $3x=12$, donc $x=4$ : \ \res{S=\ens{4}}
  \item $-3x=24$, donc $x=-8$ : \ \res{S=\ens{-8}}
  \item $3x-3x=-4-8$, soit $0=-12$, impossible : \ \res{S=\varnothing}
\end{reponses}
\rem{En d), les $x$ disparaissent et il reste une égalité fausse : aucune solution.}
\end{exo}

\begin{exo}[Équations avec des parenthèses]
\begin{reponses}
  \item $3x-12=-2x+8$, soit $5x=20$ : \ \res{S=\ens{4}}
  \item $2a-10=12-4a+8$, soit $6a=30$ : \ \res{S=\ens{5}}
  \item $m-7+2m=2m+5$, soit $3m-7=2m+5$ : \ \res{S=\ens{12}}
  \item $3p-6+4p-20=0$, soit $7p=26$ : \ \res{S=\ens{\tfrac{26}{7}}}
\end{reponses}
\rem{En c), $-(7-2m)=-7+2m$ : le $-$ devant la parenthèse change les deux signes.}
\end{exo}

\begin{exo}[Équations produit]
\begin{reponses}
  \item $3x-9=0$ ou $x+5=0$ : \ \res{S=\ens{-5\,;3}}
  \item $2x+10=0$ ou $3-4x=0$ : \ \res{S=\ens{-5\,;\tfrac{3}{4}}}
  \item $4x=0$ ou $3x-8=0$ : \ \res{S=\ens{0\,;\tfrac{8}{3}}}
  \item $2x+5=0$ : \ \res{S=\ens{-\tfrac{5}{2}}}
  \item $x(6x-5)=0$ : \ \res{S=\ens{0\,;\tfrac{5}{6}}}
  \item $(x-7)(x+7)=0$ : \ \res{S=\ens{-7\,;7}}
  \item $(x-3)\big[(2x+5)+(x-8)\big]=0$\\
        $(x-3)(3x-3)=0$ : \ \res{S=\ens{1\,;3}}
  \item $(3x+1)\big[(x+2)-(2x+6)\big]=0$\\
        $(3x+1)(-x-4)=0$ : \ \res{S=\ens{-4\,;-\tfrac{1}{3}}}
\end{reponses}
\rem{En e), diviser par $x$ ferait perdre la solution $0$ : on factorise. En d), le
carré est nul seulement si $2x+5=0$ : une seule solution.}
\end{exo}

\begin{exo}[Un premier problème]
Soit $c$ le nombre de condensateurs. Il y a $2c$ résistances et $c-8$ diodes.
\begin{calculs}
2c+c+(c-8) &= 96\\
4c-8 &= 96\\
c &= 26
\end{calculs}
Vérification : $52+26+18=96$.\par
\res{\text{52 résistances, 26 condensateurs et 18 diodes}}
\rem{« $8$ condensateurs de plus que de diodes » : ce sont les diodes qui en ont $8$ de
moins, $c-8$ et non $c+8$.}
\end{exo}

\partie{Mise en équation}

\begin{exo}[Deux commandes de câbles]
\begin{reponses}
  \item Premier club : $\res{4x}$ euros ; le second achète $\res{x+15}$ câbles et
        dépense $\res{1{,}5(x+15)}$ euros, soit $1{,}5x+22{,}5$.
  \item $4x=1{,}5x+22{,}5$, soit $2{,}5x=22{,}5$ et $x=9$.\\
        Vérification : $4\times 9=36$ et $1{,}5\times 24=36$.\\
        \res{\text{le premier club a acheté 9 câbles}}
\end{reponses}
\end{exo}

\begin{exo}[Une question d'âges]
Dans $x$ années, la technicienne aura $58+x$ ans, sa fille $33+x$ et son petit-fils
$7+x$.
\begin{calculs}
(33+x)+(7+x) &= 58+x\\
40+2x &= 58+x\\
x &= 18
\end{calculs}
Dans $18$ ans : $51+25=76$, l'âge de la technicienne. Or $76>70$ :
\res{\text{elle a tort}}.
\rem{Trouver $x$ ne suffit pas : on compare le résultat à la condition de l'énoncé
(avant $70$ ans).}
\end{exo}

\begin{exo}[Trottinettes en libre-service]
\begin{reponses}
  \item $1+0{,}25\times 12=4$ : \ \res{4\text{ €}}
  \item $1+0{,}25x=6{,}5$, soit $0{,}25x=5{,}5$ et $x=22$ : \res{22\text{ minutes}}
  \item $1+0{,}25x=0{,}35x$, soit $0{,}1x=1$ et $x=10$ : \res{10\text{ minutes}}
\end{reponses}
\end{exo}

\partie{Résoudre des inéquations}

\begin{exo}[Inéquations immédiates]
\begin{reponses}
  \item $x\le 5$ : \ \res{S=\left]-\infty\,;5\right]}
  \item $x>-3$ : \ \res{S=\left]-3\,;+\infty\right[}
  \item $x\le 5$ : \ \res{S=\left]-\infty\,;5\right]}
  \item $2x>-9$, $x>-\frac{9}{2}$ : \ \res{S=\left]-\tfrac{9}{2}\,;+\infty\right[}
  \item $x<-4$ \emph{(division par $-3$ : le sens change)} :\\
        \res{S=\left]-\infty\,;-4\right[}
  \item $-x\ge -7$, $x\le 7$ : \ \res{S=\left]-\infty\,;7\right]}
  \item $5x<2$, $x<\frac{2}{5}$ : \ \res{S=\left]-\infty\,;\tfrac{2}{5}\right[}
  \item $3-4x\ge 0$, $-4x\ge -3$, $x\le\frac{3}{4}$ :\\
        \res{S=\left]-\infty\,;\tfrac{3}{4}\right]}
\end{reponses}
\rem{Crochet fermé pour $\le$ ou $\ge$, ouvert pour $<$ ou $>$ ; toujours ouvert du
côté de l'infini.}
\end{exo}

\begin{exo}[L'inconnue des deux côtés]
\begin{reponses}
  \item $6x<8$, $x<\frac{4}{3}$ : \ \res{S=\left]-\infty\,;\tfrac{4}{3}\right[}
  \item $-3x\ge 9$, $x\le -3$ : \ \res{S=\left]-\infty\,;-3\right]}
  \item $3x-6\ge 9-2x$, $5x\ge 15$ : \ \res{S=\left[3\,;+\infty\right[}
  \item $-2x+6<4-x$, $-x<-2$, $x>2$ : \ \res{S=\left]2\,;+\infty\right[}
  \item $3x+3<5+3x$, soit $3<5$, toujours vrai : \ \res{S=\R}
  \item $2x-8\ge 2x+1$, soit $-8\ge 1$, toujours faux : \ \res{S=\varnothing}
\end{reponses}
\rem{En e) et f), les $x$ disparaissent : une inégalité vraie donne $S=\R$, une
inégalité fausse $S=\varnothing$.}
\end{exo}

\begin{exo}[Inéquations avec des fractions]
\begin{reponses}
  \item $\frac{x}{2}\le 2$, $x\le 4$ : \ \res{S=\left]-\infty\,;4\right]}
  \item $3x-2<16$, $x<6$ : \ \res{S=\left]-\infty\,;6\right[}
  \item On multiplie par $6>0$ : $3(x-1)-2(2-x)\ge 0$, soit $5x-7\ge 0$ :\\
        \res{S=\left[\tfrac{7}{5}\,;+\infty\right[}
  \item On multiplie par $4>0$ : $x-2(6-x)>12$, soit $3x>24$ :\\
        \res{S=\left]8\,;+\infty\right[}
\end{reponses}
\rem{On multiplie par un nombre \emph{positif} : le sens ne change pas. Un numérateur
précédé d'un $-$ se met entre parenthèses.}
\end{exo}

\partie{Tableaux de signes}

\begin{exo}[Signe d'une expression affine]
\noindent\petittab{$3x-12$}{$4$}{-,z,+}\hfill\petittab{$-2x+10$}{$5$}{+,z,-}\par
\smallskip
\noindent\petittab{$x+5$}{$-5$}{-,z,+}\hfill\petittab{$7-x$}{$7$}{+,z,-}\par
\smallskip
\noindent\petittab{$5x+2$}{$-\frac{2}{5}$}{-,z,+}\hfill\petittab{$-4x-6$}{$-\frac{3}{2}$}{+,z,-}
\rem{On résout d'abord $ax+b=0$ ; après la valeur trouvée, le signe est celui de $a$.}
\end{exo}

\begin{exo}[Signe d'un produit de deux facteurs]
\textbf{1)} $P(x)=(x-5)(x-2)$ ; valeurs $2$ et $5$.
\begin{tabsignes}{$x$/1,$x-2$/1,$x-5$/1,$P(x)$/1}{$-\infty$,$2$,$5$,$+\infty$}
  \tkzTabLine{,-,z,+,t,+,}
  \tkzTabLine{,-,t,-,z,+,}
  \tkzTabLine{,+,z,-,z,+,}
\end{tabsignes}
\textbf{2)} $P(x)=(1-3x)(x+4)$ ; $1-3x=0$ pour $x=\frac{1}{3}$ ; $x+4=0$ pour $x=-4$.
\begin{tabsignes}{$x$/1,$1-3x$/1,$x+4$/1,$P(x)$/1}{$-\infty$,$-4$,$\frac{1}{3}$,$+\infty$}
  \tkzTabLine{,+,t,+,z,-,}
  \tkzTabLine{,-,z,+,t,+,}
  \tkzTabLine{,-,z,+,z,-,}
\end{tabsignes}
\textbf{3)} $(2x+6)(5-x)$ ; valeurs $-3$ et $5$ ; signe $-$, $+$, $-$ : positif sur
$\left]-3\,;5\right[$, négatif ailleurs (sauf en $-3$ et $5$, où il est nul).\par
\smallskip
\textbf{4)} $x^{2}-16=(x-4)(x+4)$ ; valeurs $-4$ et $4$ ; signe $+$, $-$, $+$ :
négatif sur $\left]-4\,;4\right[$, positif sur $\left]-\infty\,;-4\right[$ et sur
$\left]4\,;+\infty\right[$.
\rem{Les valeurs de la première ligne sont rangées dans l'ordre croissant ; la dernière
ligne vient de la règle des signes, colonne par colonne.}
\end{exo}

\begin{exo}[Signe d'un produit de trois facteurs]
\textbf{1)} $P(x)=x(x-2)(x+3)$ ; valeurs $-3$, $0$ et $2$.
\begin{tabsignes}{$x$/1,$x+3$/1,$x$/1,$x-2$/1,$P(x)$/1}%
  {$-\infty$,$-3$,$0$,$2$,$+\infty$}
  \tkzTabLine{,-,z,+,t,+,t,+,}
  \tkzTabLine{,-,t,-,z,+,t,+,}
  \tkzTabLine{,-,t,-,t,-,z,+,}
  \tkzTabLine{,-,z,+,z,-,z,+,}
\end{tabsignes}
\textbf{2)} $3x(2x-5)(x+4)$ ; valeurs $-4$, $0$ et $\frac{5}{2}$ ; tous les
coefficients de $x$ sont positifs, comme en 1) : signe $-$, $+$, $-$, $+$.\par
\smallskip
\textbf{3)} $4-x^{2}=(2-x)(2+x)$ ; $P(x)=(2-x)(2+x)(x-3)$ ; valeurs $-2$, $2$ et $3$.
\begin{tabsignes}{$x$/1,$2+x$/1,$2-x$/1,$x-3$/1,$P(x)$/1}%
  {$-\infty$,$-2$,$2$,$3$,$+\infty$}
  \tkzTabLine{,-,z,+,t,+,t,+,}
  \tkzTabLine{,+,t,+,z,-,t,-,}
  \tkzTabLine{,-,t,-,t,-,z,+,}
  \tkzTabLine{,+,z,-,z,+,z,-,}
\end{tabsignes}
\rem{En 3), le facteur $2-x$ est positif \emph{avant} $2$ : le signe du produit
commence par $+$.}
\end{exo}

\partie{Inéquations produit}

\begin{exo}[Deux facteurs]
\textbf{a)} $P(x)=x(x-3)$ ; valeurs $0$ et $3$.
\begin{tabsignes}{$x$/1,$x$/1,$x-3$/1,$P(x)$/1}{$-\infty$,$0$,$3$,$+\infty$}
  \tkzTabLine{,-,z,+,t,+,}
  \tkzTabLine{,-,t,-,z,+,}
  \tkzTabLine{,+,z,-,z,+,}
\end{tabsignes}
\res{S=\left]-\infty\,;0\right]\cup\left[3\,;+\infty\right[}
\par\smallskip
Pour la suite, seule la ligne du produit est donnée.\par
\textbf{b)} Valeurs $-2$ et $6$ ; signe $+$, $-$, $+$ :
\ \res{S=\left]-2\,;6\right[}\par
\textbf{c)} Valeurs $-3$ et $5$ ; signe $-$, $+$, $-$ :
\ \res{S=\left[-3\,;5\right]}\par
\textbf{d)} Valeurs $\frac{1}{5}$ et $\frac{4}{3}$ ; signe $-$, $+$, $-$ :\par
\res{S=\left]-\infty\,;\tfrac{1}{5}\right[\cup\left]\tfrac{4}{3}\,;+\infty\right[}
\rem{Inégalité large ($\le$, $\ge$) : les valeurs qui annulent sont solutions, crochets
fermés. Inégalité stricte : crochets ouverts.}
\end{exo}

\begin{exo}[Trois facteurs]
\textbf{a)} $P(x)=x(x-2)(x+4)$ ; valeurs $-4$, $0$ et $2$.
\begin{tabsignes}{$x$/1,$x+4$/1,$x$/1,$x-2$/1,$P(x)$/1}%
  {$-\infty$,$-4$,$0$,$2$,$+\infty$}
  \tkzTabLine{,-,z,+,t,+,t,+,}
  \tkzTabLine{,-,t,-,z,+,t,+,}
  \tkzTabLine{,-,t,-,t,-,z,+,}
  \tkzTabLine{,-,z,+,z,-,z,+,}
\end{tabsignes}
\res{S=\left[-4\,;0\right]\cup\left[2\,;+\infty\right[}
\par\smallskip
\textbf{b)} Valeurs $-3$, $\frac{1}{2}$ et $2$ ; signe $+$, $-$, $+$, $-$ :\par
\res{S=\left]-3\,;\tfrac{1}{2}\right[\cup\left]2\,;+\infty\right[}\par
\textbf{c)} Valeurs $-6$, $0$ et $\frac{4}{3}$ ; signe $-$, $+$, $-$, $+$ :\par
\res{S=\left]-\infty\,;-6\right]\cup\left[0\,;\tfrac{4}{3}\right]}\par
\textbf{d)} $x(x-4)(x+4)>0$ ; valeurs $-4$, $0$ et $4$ ; signe $-$, $+$, $-$, $+$ :\par
\res{S=\left]-4\,;0\right[\cup\left]4\,;+\infty\right[}
\rem{En b), le facteur $2-x$ est positif \emph{avant} $2$ : le produit commence par
$+$.}
\end{exo}

\begin{exo}[Factoriser avant de résoudre]
\begin{reponses}
  \item $(x-6)(x+6)<0$ ; signe $+$, $-$, $+$ : \ \res{S=\left]-6\,;6\right[}
  \item $x(x-8)\le 0$ ; signe $+$, $-$, $+$ : \ \res{S=\left[0\,;8\right]}
  \item $(3x-2)(3x+2)(x+2)>0$ ; valeurs $-2$, $-\frac{2}{3}$, $\frac{2}{3}$ ; signe
        $-$, $+$, $-$, $+$ :\\
        \res{S=\left]-2\,;-\tfrac{2}{3}\right[\cup\left]\tfrac{2}{3}\,;+\infty\right[}
  \item $\big[(x+2)-4\big]\big[(x+2)+4\big]>0$, soit $(x-2)(x+6)>0$ :\\
        \res{S=\left]-\infty\,;-6\right[\cup\left]2\,;+\infty\right[}
\end{reponses}
\rem{$x^{2}<36$ ne donne pas seulement $x<6$ : le tableau de signes évite cette
erreur.}
\end{exo}

\partie{Inéquations quotient}

\begin{exo}[Premières inéquations quotient]
\textbf{1)} $x+5=0$ pour $x=\res{-5}$ ; \ $2-x=0$ pour $x=\res{2}$ ; \
$3x-9=0$ pour $x=\res{3}$.\par
\smallskip
\textbf{2)} \textbf{a)} Valeur interdite $-5$ ; le numérateur s'annule en $4$.
\begin{tabsignes}{$x$/1,$4-x$/1,$x+5$/1,$\frac{4-x}{x+5}$/1.2}{$-\infty$,$-5$,$4$,$+\infty$}
  \tkzTabLine{,+,t,+,z,-,}
  \tkzTabLine{,-,z,+,t,+,}
  \tkzTabLine{,-,d,+,z,-,}
\end{tabsignes}
\res{S=\left]-5\,;4\right[}\par
\smallskip
\textbf{b)} Interdite $3$ ; numérateur nul en $-2$ ; signe $+$, $-$, $+$ :\par
\res{S=\left[-2\,;3\right[}\par
\textbf{c)} Interdite $2$ ; numérateur nul en $\frac{7}{2}$ ; signe $+$, $-$, $+$ :\par
\res{S=\left]-\infty\,;2\right[\cup\left[\tfrac{7}{2}\,;+\infty\right[}\par
\textbf{d)} Interdite $4$ ; numérateur nul en $-\frac{2}{3}$ ; signe $-$, $+$, $-$ :\par
\res{S=\left]-\infty\,;-\tfrac{2}{3}\right[\cup\left]4\,;+\infty\right[}
\rem{Même avec $\le$ ou $\ge$, la valeur interdite a toujours un crochet ouvert.}
\end{exo}

\begin{exo}[Avec une factorisation]
\textbf{a)} $Q(x)=\dfrac{x(x-2)}{5-2x}$ ; interdite $\frac{5}{2}$ ; numérateur nul en
$0$ et $2$.
\begin{tabsignes}{$x$/1,$x$/1,$x-2$/1,$5-2x$/1,$Q(x)$/1}%
  {$-\infty$,$0$,$2$,$\frac{5}{2}$,$+\infty$}
  \tkzTabLine{,-,z,+,t,+,t,+,}
  \tkzTabLine{,-,t,-,z,+,t,+,}
  \tkzTabLine{,+,t,+,t,+,z,-,}
  \tkzTabLine{,+,z,-,z,+,d,-,}
\end{tabsignes}
\res{S=\left[0\,;2\right]\cup\left]\tfrac{5}{2}\,;+\infty\right[}\par
\smallskip
\textbf{b)} $\dfrac{(x-4)(x+4)}{2-x}>0$ ; interdite $2$ ; numérateur nul en $-4$ et
$4$ ; signe $+$, $-$, $+$, $-$ :\par
\res{S=\left]-\infty\,;-4\right[\cup\left]2\,;4\right[}\par
\textbf{c)} $\dfrac{6-3x}{(x-3)(x+3)}\le 0$ ; interdites $-3$ et $3$ ; numérateur nul
en $2$ ; signe $+$, $-$, $+$, $-$ :\par
\res{S=\left]-3\,;2\right]\cup\left]3\,;+\infty\right[}
\end{exo}

\begin{exo}[Tout ramener d'un même côté]
\textbf{a)} Interdite $-3$.
\begin{calculs}
\frac{3x+1}{x+3} &\le 2\\
\frac{3x+1-2(x+3)}{x+3} &\le 0\\
\frac{x-5}{x+3} &\le 0
\end{calculs}
Numérateur nul en $5$ ; signe $+$, $-$, $+$ : \ \res{S=\left]-3\,;5\right]}\par
\textbf{b)} Interdite $2$.
\begin{calculs}
\frac{2-5x}{2-x} &\ge 3\\
\frac{2-5x-3(2-x)}{2-x} &\ge 0\\
\frac{-2x-4}{2-x} &\ge 0
\end{calculs}
Numérateur nul en $-2$ ; $-2x-4$ et $2-x$ sont positifs avant leur valeur ; signe
$+$, $-$, $+$ :\par
\res{S=\left]-\infty\,;-2\right]\cup\left]2\,;+\infty\right[}\par
\textbf{c)} Interdite $\frac{3}{2}$.
\begin{calculs}
\frac{x+4}{3-2x} &< 1\\
\frac{x+4-(3-2x)}{3-2x} &< 0\\
\frac{3x+1}{3-2x} &< 0
\end{calculs}
Numérateur nul en $-\frac{1}{3}$ ; signe $-$, $+$, $-$ :\par
\res{S=\left]-\infty\,;-\tfrac{1}{3}\right[\cup\left]\tfrac{3}{2}\,;+\infty\right[}
\rem{Multiplier par $x+3$ serait faux : son signe dépend de $x$, on ne saurait pas
s'il faut changer le sens.}
\end{exo}

\partie{Problèmes avec des inéquations}

\begin{exo}[Seuil de rentabilité]
\begin{questions}
  \item $C(25)=15\times 25+480=855$ et $R(25)=31\times 25=775$.\\
        $R(25)<C(25)$ : \res{\text{pas de bénéfice}}.
  \item $31x>15x+480$, soit $16x>480$ et $x>30$.\\
        \res{\text{bénéfice à partir de 31 modules par jour}}
  \item $\dfrac{15x+480}{x}\le 21$, soit $\dfrac{480-6x}{x}\le 0$.\\
        $x>0$ : le quotient a le signe de $480-6x$.\\
        $480-6x\le 0$, soit $x\ge 80$ : \res{\text{au moins 80 modules par jour}}
\end{questions}
\rem{$x$ est un nombre entier de modules : au-dessus de $30$, le premier entier est
$31$.}
\end{exo}

\begin{exo}[Louer un utilitaire électrique]
\begin{questions}
  \item $60$~km : A coûte $30$~€, B coûte $36+12=48$~€ : \res{\text{A}}.\\
        $150$~km : A coûte $75$~€, B coûte $36+30=66$~€ : \res{\text{B}}.
  \item \res{A(x)=0{,}5x} \quad et \quad \res{B(x)=36+0{,}2x}
  \item $36+0{,}2x<0{,}5x$, soit $36<0{,}3x$ et $x>120$.\\
        \res{\text{B est moins cher au-delà de 120 km}}
  \item $36+0{,}2x\le 80$, soit $0{,}2x\le 44$ et $x\le 220$ :
        \res{\text{220 km au plus}}
  \item $C(x)=60+0{,}1x$.\\
        $C(x)<A(x)$ : $60<0{,}4x$, soit $x>150$.\\
        $C(x)<B(x)$ : $24<0{,}1x$, soit $x>240$.\\
        Les deux à la fois, avec $x\le 300$ : \res{x\in\left]240\,;300\right]}
\end{questions}
\rem{« Moins cher que les deux autres » : les deux conditions sont vraies en même
temps ; on garde la plus exigeante.}
\end{exo}

\begin{exo}[Bénéfice d'un fabricant de trottinettes]
\begin{questions}
  \item $B(10)=-200+2\,400-4\,000$, soit \res{-1\,800} : perte.\\
        $B(50)=-5\,000+12\,000-4\,000$, soit \res{3\,000} : bénéfice.
  \item \begin{calculs}
        (2x-40)(100-x)&=200x-2x^{2}-4\,000+40x\\
        (2x-40)(100-x)&=-2x^{2}+240x-4\,000
        \end{calculs}
        Donc \res{B(x)=(2x-40)(100-x)}.
  \item Valeurs $20$ et $100$.
\begin{tabsignes}{$x$/1,$2x-40$/1,$100-x$/1,$B(x)$/1}{$0$,$20$,$100$,$120$}
  \tkzTabLine{,-,z,+,t,+,}
  \tkzTabLine{,+,t,+,z,-,}
  \tkzTabLine{,-,z,+,z,-,}
\end{tabsignes}
        \res{S=\left]20\,;100\right[}
  \item \res{\text{de 21 à 99 trottinettes par semaine}}
  \item \begin{calculs}
        -2(x-40)(x-80)&=-2\left(x^{2}-120x+3\,200\right)\\
        -2(x-40)(x-80)&=-2x^{2}+240x-6\,400
        \end{calculs}
        C'est bien $B(x)-2\,400$.\\
        $B(x)\ge 2\,400$ s'écrit $-2(x-40)(x-80)\ge 0$ ; on divise par $-2<0$, le sens
        change : $(x-40)(x-80)\le 0$, signe $+$, $-$, $+$.\\
        \res{\text{de 40 à 80 trottinettes par semaine}}
\end{questions}
\rem{Le tableau est fait sur $\intff{0}{120}$ : la première colonne commence à $0$,
pas à $-\infty$.}
\end{exo}

\end{document}
